\chapter{Method}
\label{ch:method}
Describe your approach in enough detail to be reproducible. Inline maths
such as $f(x)=\sigma(\mathbf{w}^\top\mathbf{x}+b)$ sits naturally in the
Libertinus body, and display equations are numbered per chapter:
\begin{equation}
  \mathcal{L}(\theta)=\frac{1}{N}\sum_{i=1}^{N}\ell\!\left(f_\theta(x_i),y_i\right).
  \label{eq:loss}
\end{equation}

\begin{figure}[h!]
  \includegraphics[width=0.50\linewidth, center]{adafruit-i2s-mems-microphone-breakout.pdf}
  \caption{Adafruit MEMS board}
  \label{fig:abc}
\end{figure}

Figure~\ref{fig:signal} is loaded with \texttt{\textbackslash includegraphics\{example-plot\}};
no path is needed because the style file sets \texttt{\textbackslash graphicspath}.
\begin{figure}[t]
  \centering
  \includegraphics[width=0.72\textwidth]{example-plot}
  \caption{A damped signal with observations, stored under \texttt{figures/}.}
  \label{fig:signal}
\end{figure}

\section{A table}
\begin{table}[t]
  \centering
  \caption{Results against two baselines.}
  \label{tab:results}
  \begin{tabular}{lcc}
    \toprule
    \textbf{Configuration} & \textbf{Accuracy} & \textbf{Time (s)}\\
    \midrule
    Baseline A & 0.81 & 12.4\\
    Baseline B & 0.86 & \phantom{0}9.7\\
    \textbf{Ours} & \textbf{0.92} & \phantom{0}8.1\\
    \bottomrule
  \end{tabular}
\end{table}
Reference it as Table~\ref{tab:results} and Equation~\eqref{eq:loss}.
